\documentclass[10pt,a4paper]{article}
\usepackage[margin=2.2cm]{geometry}
\usepackage{fontspec}
\setmonofont{DejaVu Sans Mono}[Scale=0.84]
\newfontfamily\unifont{DejaVu Serif}[Scale=0.9]
\usepackage{amsmath,amssymb,mathtools}
\setcounter{MaxMatrixCols}{20}
\usepackage{longtable,booktabs,array}
\usepackage{fancyvrb}
\usepackage[dvipsnames]{xcolor}
\usepackage[colorlinks=true,linkcolor=NavyBlue,urlcolor=NavyBlue]{hyperref}
\usepackage{microtype}
\usepackage{needspace}
\usepackage{tocloft}\setlength{\cftbeforesecskip}{0.45em}
\setlength{\parindent}{0pt}\setlength{\parskip}{0.55em}
\setlength{\emergencystretch}{8em}
\makeatletter
\renewcommand\section{\@startsection{section}{1}{\z@}{-3.5ex \@plus -1ex \@minus -.2ex}{2.3ex \@plus.2ex}{\normalfont\Large\bfseries\raggedright}}
\renewcommand\subsection{\@startsection{subsection}{2}{\z@}{-3.25ex\@plus -1ex \@minus -.2ex}{1.5ex \@plus .2ex}{\normalfont\large\bfseries\raggedright}}
\renewcommand\subsubsection{\@startsection{subsubsection}{3}{\z@}{-3.25ex\@plus -1ex \@minus -.2ex}{1.5ex \@plus .2ex}{\normalfont\normalsize\bfseries\raggedright}}
\makeatother
\renewcommand{\arraystretch}{1.15}
\allowdisplaybreaks
\begin{document}
\title{The Standard Model on the world brane of 7+1 spacetime: provenance\\[0.4em]\large Sources, data, decisions, derivations and the verification record of the Standard Model on the brane, the portals of dirac16complex and neutrino flavour conversion — sections 17–19 of the notebook 7+1.nb}
\author{Repository: github.com/once-ere/7plus1 — prepared for Nash\\ Notebook: 7+1.nb (Wolfram Language 15.0.1)\\ Document generated by tools/make\_docs.py on 2026-10-28}
\date{}
\maketitle
\begin{abstract}
This record states what was asked, which experimental data were used and where they come from, every decision taken with the equation behind it, how each result was derived and how it was checked: the 24 checks of sections 17–19 of 7+1.nb (24 True), the 33 boolean checks of the independent Python program tools/ref\_sm.py (33 True), the comparison of every number printed by both programs, and the SHA-256 hashes of the files.
\end{abstract}
\tableofcontents
\Needspace{12\baselineskip}\section*{1. The task}\addcontentsline{toc}{section}{1. The task}
This stage of the project asked (summarized; the original wording is the author’s private working material) for the TOTALITY of the Standard Model together with gravity and the field dirac16complex on the 7+1 spacetime of this project (one time coordinate $x_0$, seven space coordinates $x_1,\dots,x_7$), following the ideas of the repository github.com/once-ere/Dirac\_claude — there the Standard Model lives on a seven-dimensional fold of a 4+4 spacetime — but not its metric. Here the Standard Model lives on the 3+1-dimensional world brane $x_4=\dots=x_7=0$, the auxiliary coordinates form a torus, and gravity and dirac16complex live in the whole spacetime. Every documentation file has a companion provenance file and a symbol file, in Markdown and LaTeX with PDFs; no document refers the reader to another one, and every equation used is written in the document that uses it. On 27 October 2026 the author added: make clear in all documentation, examples and provenance records that the xAct list $\{7,1,0\}$ declares seven space-like dimensions and one time-like dimension, while $\mathrm{diag}(+1,-1,\dots,-1)$ is simply the squares of the octonion basis (section 1.1).

\Needspace{8\baselineskip}\subsection*{1.1 The signature: the author’s explanation of 27 October 2026}\addcontentsline{toc}{subsection}{1.1 The signature: the author’s explanation of 27 October 2026}
\textbf{Source.} On 27 October 2026 the author of the project explained that the xAct list $\{7,1,0\}$ declares seven space-like dimensions and one time-like dimension, while $\mathrm{diag}(+1,-1,\dots,-1)$ is simply the squares of the octonion basis, and asked that this be made clear in all documentation, examples and provenance records. \textbf{What it corrected.} The version of 7+1.nb pushed in the commits 52d8729 to dc2f178 of github.com/once-ere/7plus1 declared the xAct metric of section 12 with \texttt{DefMetric[\{1, 7, 0\}, …]}, reading the list as the signs of $\eta$; since then section 12 declares \texttt{DefMetric[\{7, 1, 0\}, g71[-a, -b], CD, …]}, exactly as the author’s notebook Einstein-Lovelock-7+1-Nash.nb does. \textbf{What was added.} Checks in 7+1.nb — section 4.7: $\eta$ is the squares of the octonion basis, $e_Ae_A=\eta_{AA}e_0$ and $\mathrm{Re}(e_Ae_B)=\eta_{AB}$; section 12.2: \texttt{SignatureOfMetric[g71]} $=\{7,1,0\}$ with \texttt{SignDetOfMetric[g71]} $=-1$, $\mathrm{Re}(e_Ae_B)=\eta_{AB}$, and $\mathrm{Re}(x\,x)=\eta_{AB}x^Ax^B$; section 17: the flat metric $\mathrm{diag}(+1,-1,-1,-1)$ of the world brane is the squares of the basis $e_0,e_1,e_2,e_3$ of the quaternion subalgebra — and the same statements in the Python programs \texttt{tools/\allowbreak{}ref71.py}, \texttt{tools/\allowbreak{}ref\_\allowbreak{}gravity.py} and \texttt{tools/\allowbreak{}ref\_\allowbreak{}sm.py}; the explanation itself in the text of 7+1.nb (sections 1.2, 1.4, 4.2, 4.7, 7, 11, 12.2, 13, 15.1, 17) and in every document, symbol file and provenance file of the repository. \textbf{What did not change.} No result. Compared output by output with the backup below, every output of 7+1.nb is the same, except the declared list itself, the new checks and the check counts (and, as in every run, the name of the private context and the measured times); in particular every component of every Lovelock tensor is unchanged. The components follow from the component matrix of the metric, and the sign of the determinant, \texttt{SignDetOfMetric}, is $-1$ for both lists. \textbf{Backup.} The evaluated notebook as it was before is kept in the workspace as \texttt{backups/\allowbreak{}7+1\_\allowbreak{}2026-10-27\_\allowbreak{}before\_\allowbreak{}signature.nb} (the folder backups/ is excluded from the repository by .gitignore), 1034401 bytes, SHA-256 \texttt{06094248\allowbreak{}8caf4f45\allowbreak{}120ffd37\allowbreak{}be0123e9\allowbreak{}6992a31c\allowbreak{}d7fdbde0\allowbreak{}9ee4dc4c\allowbreak{}ed4f5d3a}; \texttt{tools/\allowbreak{}verify\_\allowbreak{}7p1\_\allowbreak{}nb.py} compares every input, text and section cell of 7+1.nb with it and finds every cell present, each one unchanged or changed by this explanation, and it compares every output in the same way.

\textbf{The explanation, in full.} The 7+1 spacetime has one time-like dimension ($x_0$, the direction of the octonion unit $e_0=1$) and seven space-like dimensions ($x_1,x_2,x_3$ ordinary space, $x_4,\dots,x_7$ auxiliary space). Two different objects describe this, and neither may be replaced by the other. (i) The tensor package xAct, used in section 12 of 7+1.nb exactly as in the author’s notebook Einstein-Lovelock-7+1-Nash.nb, declares the metric with \texttt{DefMetric[\{7, 1, 0\}, g71[-a, -b], CD, …]}: the list $\{7,1,0\}$ declares seven space-like dimensions, one time-like dimension and no degenerate dimension — the numbers of entries $+1$ and $-1$ of $\eta_{71}=\mathrm{diag}(-1,+1,\dots,+1)$ (\texttt{η71}), the same spacetime written in the sign convention of 4+4.nb (space-like $=+1$, time-like $=-1$). (ii) The matrix $\eta=-\eta_{71}=\mathrm{diag}(+1,-1,\dots,-1)$ (\texttt{η17}) is simply the squares of the octonion basis: $e_0^2=+1$ for the unit $e_0=1$ (time) and $e_i^2=-1$ for the seven imaginary units $e_1,\dots,e_7$ (space); more generally it is the real part of the product of two basis units, and it gives the quadratic form of the octonions:
\begin{gather*}
\eta_{AB}=\mathrm{Re}(e_Ae_B),\qquad e_0^2=+1,\qquad e_i^2=-1\ \ (i=1,\dots,7),\\ \mathrm{Re}(x\,x)=x_0^2-x_1^2-\cdots-x_7^2=\eta_{AB}\,x^Ax^B\qquad (x=x^Ae_A\in\mathbb O).
\end{gather*}
$\eta$ gives the signs of the Clifford relation $\Gamma_A\Gamma_B+\Gamma_B\Gamma_A=2\,\eta_{AB}\,1_{16}$ of the real 16×16 Dirac matrices: time squares to $+1_{16}$ and each of the seven space directions to $-1_{16}$, exactly like the octonion units (in the notation $Cl(p,q)$, with $p$ generators squaring to $+1$ and $q$ to $-1$, the algebra $Cl(1,7)\cong\mathbb R(16)$). It also gives the signs of the metric components: $g_{00}>0$ along the time-like direction, negative squares along the space-like directions. Every component computed by 7+1.nb follows from the component matrix of the metric; the declared list fixes for xAct the sign of the determinant of the metric, \texttt{SignDetOfMetric[g71]} $=-1$, which is the same in both sign conventions, since $\det \eta=\det \eta_{71}=-1$. 7+1.nb checks \texttt{SignatureOfMetric[g71]} $=\{7,1,0\}$, \texttt{SignDetOfMetric[g71]} $=-1$, $\mathrm{Re}(e_Ae_B)=\eta_{AB}$ and $\mathrm{Re}(x\,x)=\eta_{AB}x^Ax^B$.

On the world brane the same reading holds: one time-like and three space-like dimensions (in the counting of xAct the list $\{3,1,0\}$), and the flat frame metric $\mathrm{diag}(+1,-1,-1,-1)$ is $\eta$ restricted to $x_0,\dots,x_3$, the squares of the basis $e_0=1,e_1,e_2,e_3$ of the quaternion subalgebra $\mathbb H$ of the octonions (a check of section 17 of 7+1.nb and of \texttt{tools/\allowbreak{}ref\_\allowbreak{}sm.py}).

\Needspace{12\baselineskip}\section*{2. Data and their sources}\addcontentsline{toc}{section}{2. Data and their sources}
{\footnotesize
\begin{longtable}{>{\raggedright\arraybackslash}p{\dimexpr 0.3000\linewidth-2\tabcolsep\relax}>{\raggedright\arraybackslash}p{\dimexpr 0.3800\linewidth-2\tabcolsep\relax}>{\raggedright\arraybackslash}p{\dimexpr 0.3200\linewidth-2\tabcolsep\relax}}
\toprule
\textbf{quantity} & \textbf{value used} & \textbf{source} \\
\midrule\endfirsthead
\toprule
\textbf{quantity} & \textbf{value used} & \textbf{source} \\
\midrule\endhead
\endfoot
\bottomrule\endlastfoot
$G_F$ & $1.1663788\times10^{-5}$ GeV$^{-2}$ & PDG 2025 \\
$m_W$, $m_Z$, $m_H$ & 80.3692, 91.1880, 125.20 GeV & PDG 2025 \\
CKM: $s_{12}$, $s_{13}$, $s_{23}$, $\delta$ & 0.22501, 0.003732, 0.04183, 1.147 & PDG 2024 (CKM review, global fit) \\
$J$, $A$, $\bar\rho$, $\bar\eta$ (for comparison) & $(3.12^{+0.13}_{-0.12})\times10^{-5}$; $0.826\pm0.016$; $0.1591\pm0.0094$; $0.3523\pm0.0073$ & PDG 2024 \\
$m_u$, $m_d$, $m_s$ & 2.16, 4.70, 93.5 MeV ($\overline{\rm MS}$, 2 GeV) & PDG 2025 \\
$m_c$, $m_b$, $m_t$ & 1.2730 GeV, 4.183 GeV ($\overline{\rm MS}$), 172.56 GeV (direct) & PDG 2025 \\
$\sin^2\theta_{12}$, $\sin^2\theta_{23}$, $\sin^2\theta_{13}$, $\delta$ & 0.307, 0.561, 0.02195, 177° & NuFit-6.0, normal ordering, without SK atmospheric data \\
$\Delta m^2_{21}$, $\Delta m^2_{31}$ & $7.49\times10^{-5}$, $2.534\times10^{-3}$ eV² & NuFit-6.0 (same set) \\
$H_0$, $\Omega_m$, $\Omega_r$ & 67.4 km/s/Mpc, 0.315, $9.14\times10^{-5}$ & Planck 2018 ($t_0=13.787\pm0.020$ Gyr for comparison) \\
$M_{\rm Pl}$ (reduced) & $2.435\times10^{18}$ GeV & standard value \\
$\hbar c$ & $1.973269804\times10^{-16}$ GeV m & CODATA \\
$N_A$ & $6.02214076\times10^{23}$ mol$^{-1}$ & SI (exact) \\
$T_\gamma$ & 2.7255 K & CMB temperature (Fixsen 2009) \\
$k_B$ & $8.617333262\times10^{-5}$ eV/K & SI (exact) \\
$R_\odot$; $n_e(r)$ & $6.957\times10^8$ m; $245N_Ae^{-10.54r/R_\odot}$ cm$^{-3}$ & IAU nominal radius; exponential fit to the standard solar model (Bahcall 1989) \\
DUNE: $L$, $E$, $\rho$, $Y_e$ & 1300 km, 2.5 GeV, 2.848 g/cm³, 0.5 & DUNE design baseline and crust density \\*
\end{longtable}}
Full references: S. Navas et al. (Particle Data Group), Phys. Rev. D 110 (2024) 030001 and the 2025 update (pdg.lbl.gov); I. Esteban, M. C. Gonzalez-Garcia, M. Maltoni, I. Martinez-Soler, J. P. Pinheiro, T. Schwetz, JHEP 12 (2024) 216 (NuFit-6.0); Planck Collaboration, Astron. Astrophys. 641 (2020) A6; D. J. Fixsen, Astrophys. J. 707 (2009) 916; J. N. Bahcall, \emph{Neutrino Astrophysics} (Cambridge, 1989).

Software: Wolfram Language 15.0.1 (front end WolframNB.exe); Python 3.13.17 with numpy 2.5.3, scipy 1.18.1, sympy 1.14.0.

\Needspace{12\baselineskip}\section*{3. Decisions}\addcontentsline{toc}{section}{3. Decisions}
{\footnotesize
\begin{longtable}{>{\raggedright\arraybackslash}p{\dimexpr 0.0550\linewidth-2\tabcolsep\relax}>{\raggedright\arraybackslash}p{\dimexpr 0.4750\linewidth-2\tabcolsep\relax}>{\raggedright\arraybackslash}p{\dimexpr 0.4700\linewidth-2\tabcolsep\relax}}
\toprule
\textbf{\#} & \textbf{decision} & \textbf{reason / equation} \\
\midrule\endfirsthead
\toprule
\textbf{\#} & \textbf{decision} & \textbf{reason / equation} \\
\midrule\endhead
\endfoot
\bottomrule\endlastfoot
1 & the Standard Model lives on the brane $W=\lbrace y=0\rbrace$ of a torus $T^4$ of the auxiliary coordinates & the 7+1 analogue of the fold of Dirac\_claude; $\int_Wd^4x\sqrt{|h|}=\int d^8X\sqrt{|g|}\,\delta^4(y)/b^4$ \\
2 & three generations with right-handed neutrinos (Dirac neutrinos) & anomalies cancel with or without $\nu_R$; $\nu_R$ gives Dirac masses $m=yv/\sqrt2$ \\
3 & tree-level electroweak relations with the on-shell weak angle & $\cos\theta_W=m_W/m_Z$ makes $\rho=m_W^2/(m_Z^2\cos^2\theta_W)=1$ exactly \\
4 & standard parametrization of both mixing matrices & $U=R_{23}U_{13}(\delta)R_{12}$; Yukawa matrices in the basis where the up-type (charged-lepton) Yukawa is diagonal \\
5 & lightest neutrino massless, normal ordering & the oscillation data fix only $\Delta m^2_{ij}$ \\
6 & brane decomposition with the Dirac representation and $\rho_4=i\gamma_0$, $\rho_k=-\gamma_{k-4}$ ($k=5,6,7$) & gives anti-Hermitian $\rho_k$ with $\lbrace\rho_k,\rho_l\rbrace=-2\delta_{kl}$, so $R_k=\gamma_5\otimes\rho_k$ squares to $-1$ \\
7 & the intertwiner by a Schur average & $\sum_K\Gamma^{(K)}s(R^{(K)})^{-1}$ commutes in the sense $\Gamma_AU=UR_A$ for any $s$; non-zero for a suitable matrix unit $s$ \\
8 & portals are invented couplings at the brane, separate for the two chiral halves & as in Dirac\_claude; they are the only contact between dirac16complex and the Standard Model besides gravity \\
9 & neutrino evolution in matter by the exponential of the constant-density Hamiltonian; the Sun by direct integration of the two-flavour equation & exact for constant density; checked against the adiabatic formula \\
10 & the solar two-flavour Hamiltonian uses the effective potential $c_{13}^2V_{CC}$ and $P_{ee}=c_{13}^4P^{(2)}_{ee}+s_{13}^4$ & standard reduction of three flavours when $\Delta m^2_{31}\gg2EV$ \\*
\end{longtable}}
\Needspace{12\baselineskip}\section*{4. Derivations}\addcontentsline{toc}{section}{4. Derivations}
\textbf{Anomaly sums} (one generation, left-handed $+$, right-handed $-$):
\begin{align*}
A_{SU(3)^2U(1)}&=2\cdot\tfrac16-\tfrac23+\tfrac13=0,\qquad A_{SU(2)^2U(1)}=3\cdot\tfrac16-\tfrac12=0,\qquad A_{SU(3)^3}=2-1-1=0, \\
A_{U(1)^3}&=6\big(\tfrac16\big)^3-3\big(\tfrac23\big)^3-3\big(-\tfrac13\big)^3+2\big(-\tfrac12\big)^3-(-1)^3=\tfrac1{36}-\tfrac89+\tfrac19-\tfrac14+1=0, \\
A_{\rm grav^2U(1)}&=6\cdot\tfrac16-3\cdot\tfrac23-3\cdot\big(-\tfrac13\big)+2\cdot\big(-\tfrac12\big)-(-1)=0 .
\end{align*}
\textbf{Gauge-boson masses.} With $\langle\Phi\rangle=(0,v/\sqrt2)$:
\begin{equation*}
\big|D_\mu\langle\Phi\rangle\big|^2=\frac{v^2}8\Big[g^2\big((W^1)^2+(W^2)^2\big)+\big(gW^3-g'B\big)^2\Big],
\end{equation*}
whose Hessian in $(W^1,W^2,W^3,B)$ has the eigenvalues $g^2v^2/4$ (twice), $(g^2+g'^2)v^2/4$ and $0$.

\textbf{Brane bilinears.} $\Gamma_AU=UR_A$ and $U^\dagger U=1$ give $\bar\psi\Gamma^\mu\psi'=\psi^\dagger\Gamma_0\Gamma^\mu\psi'=\zeta^\dagger R_0R^\mu\zeta'$, and $R_0R^\mu=(\gamma_0\gamma^\mu)\otimes1_4$ for $\mu=0,\dots,3$, so the transverse index $a$ is summed: $\sum_a\bar\psi_a\gamma^\mu\psi'_a$.

\textbf{Kaluza–Klein masses.} $X=\gamma_5\otimes\vec\rho\cdot\vec q/b_0$ anticommutes with $\gamma_\mu\otimes1_4$ ($\gamma_5$ anticommutes with $\gamma_\mu$) and $X^2=\gamma_5^2\otimes(\vec\rho\cdot\vec q)^2/b_0^2=-|\vec q|^2/b_0^2$; then $(\gamma\cdot p\otimes1+X)^2=(p^2-|\vec q|^2/b_0^2)\,1$, so the operator $\gamma\cdot p\otimes1+X-m$ is singular exactly when $p^2=m^2+|\vec q|^2/b_0^2$.

\textbf{Non-unitary mixing.} For $\mathcal M=\begin{pmatrix}m_D&\mu\\0&m1_4\end{pmatrix}$ with $m\gg\mu,m_D$ the left singular vectors give, to second order in $\mu/m$, $N=(1-\tfrac12\mu\mu^\dagger/m^2)U$, hence $1-NN^\dagger=\mu\mu^\dagger/m^2$.

\textbf{Vacuum oscillations.} Expanding $|\sum_iU_{\beta i}U^*_{\alpha i}e^{-im_i^2L/(2E)}|^2$ and using unitarity, $\sum_iU_{\beta i}U^*_{\alpha i}=\delta_{\alpha\beta}$:
\begin{equation*}
P(\alpha\to\beta)=\delta_{\alpha\beta}-4\sum_{i>j}\mathrm{Re}\,X_{ij}\sin^2\Delta_{ij}+2\sum_{i>j}\mathrm{Im}\,X_{ij}\sin2\Delta_{ij},\qquad X_{ij}=U^*_{\alpha i}U_{\beta i}U_{\alpha j}U^*_{\beta j}.
\end{equation*}
(An earlier version of the notebook used $X_{ij}=U_{\alpha i}U^*_{\beta i}U^*_{\alpha j}U_{\beta j}$, the complex conjugate, which flips the sign of the CP-odd term; it was corrected, and the check compares the formula with the direct evolution for all nine $(\alpha,\beta)$.)

\textbf{MSW angle.} Diagonalizing the two-flavour Hamiltonian in matter gives the mixing angle in matter $\theta_m$:
\begin{equation*}
\frac{\Delta m^2}{4E}\begin{pmatrix}-\cos2\theta&\sin2\theta\\\sin2\theta&\cos2\theta\end{pmatrix}+\mathrm{diag}(V,0)\quad\Longrightarrow\quad \tan2\theta_m=\frac{\sin2\theta}{\cos2\theta-2EV/\Delta m^2} .
\end{equation*}
\Needspace{12\baselineskip}\section*{5. Verification}\addcontentsline{toc}{section}{5. Verification}
\textbf{Evaluation of 7+1.nb}: Wolfram 15.0.1 for Microsoft Windows (64-bit) (July 2, 2026); finished Wed 28 Oct 2026 08:57:47; message cells 0. The checks of sections 17–19:

{\footnotesize
\begin{longtable}{>{\raggedright\arraybackslash}p{\dimexpr 0.0550\linewidth-2\tabcolsep\relax}>{\raggedright\arraybackslash}p{\dimexpr 0.7450\linewidth-2\tabcolsep\relax}>{\raggedright\arraybackslash}p{\dimexpr 0.0800\linewidth-2\tabcolsep\relax}>{\raggedright\arraybackslash}p{\dimexpr 0.1200\linewidth-2\tabcolsep\relax}}
\toprule
\textbf{\#} & \textbf{check} & \textbf{result} & \textbf{section} \\
\midrule\endfirsthead
\toprule
\textbf{\#} & \textbf{check} & \textbf{result} & \textbf{section} \\
\midrule\endhead
\endfoot
\bottomrule\endlastfoot
1 & \texttt{brane frame metric diag(+1, -1, -1, -1) == eta17 restricted to x0..x3: the squares of the octonion basis e[0], e[1], e[2], e[3] of the quaternion subalgebra (Re(e[A] e[B]), A, B = 0..3), which is closed under the octonion product} & True & 17. \\
2 & \texttt{anomaly cancellation per generation (with nu\_\allowbreak{}R): SU(3)\textasciicircum{}2U(1), SU(2)\textasciicircum{}2U(1), U(1)\textasciicircum{}3, grav\textasciicircum{}2U(1), SU(3)\textasciicircum{}3 all vanish; doublets per generation even} & True & 17.1 \\
3 & \texttt{electric charges Q = T3 + Y: (u, d) = (2/\allowbreak{}3, -1/\allowbreak{}3), (nu, e) = (0, -1), u\_\allowbreak{}R 2/\allowbreak{}3, d\_\allowbreak{}R -1/\allowbreak{}3, e\_\allowbreak{}R -1, nu\_\allowbreak{}R 0} & True & 17.1 \\
4 & \texttt{g* = 28 + (7/\allowbreak{}8) 90 = 106.75 for the Standard Model fields (without nu\_\allowbreak{}R)} & True & 17.1 \\
5 & \texttt{gauge-boson mass matrix: eigenvalues g\textasciicircum{}2 v\textasciicircum{}2/\allowbreak{}4 (twice), (g\textasciicircum{}2 + g'\textasciicircum{}2) v\textasciicircum{}2/\allowbreak{}4 and 0 (photon)} & True & 17.3 \\
6 & \texttt{tree level: m\_\allowbreak{}W == g v/\allowbreak{}2, m\_\allowbreak{}Z == Sqrt[g\textasciicircum{}2 + g'\textasciicircum{}2] v/\allowbreak{}2, rho == m\_\allowbreak{}W\textasciicircum{}2/\allowbreak{}(m\_\allowbreak{}Z\textasciicircum{}2 cos\textasciicircum{}2 theta\_\allowbreak{}W) == 1, m\_\allowbreak{}H\textasciicircum{}2 == 2 lambda v\textasciicircum{}2} & True & 17.3 \\
7 & \texttt{CKM: unitary; Jarlskog invariant J = 3.12 x 10\textasciicircum{}-5 (PDG 2024: 3.12 +0.13 -0.12 x 10\textasciicircum{}-5); A, rhobar, etabar within the PDG uncertainties} & True & 17.4 \\
8 & \texttt{Yukawa matrices: singular values x v/\allowbreak{}Sqrt[2] are the quark masses and |U\_\allowbreak{}uL\textasciicircum{}dag.U\_\allowbreak{}dL| == |V\_\allowbreak{}CKM|} & True & 17.4 \\
9 & \texttt{PMNS unitary; Dirac neutrino Yukawa Sqrt[2] U.diag(m)/\allowbreak{}v reproduces the masses (singular values)} & True & 17.4 \\
10 & \texttt{age of the universe from the brane Friedmann equation: 13.79 Gyr (Planck 2018: 13.787 +- 0.020 Gyr)} & True & 17.5 \\
11 & \texttt{ADD relation: kappa8 rho4/\allowbreak{}V\_\allowbreak{}perp == 8 Pi G4 rho4 with G4 = G8/\allowbreak{}V\_\allowbreak{}perp (kappa8 = 8 Pi G8)} & True & 17.5 \\
12 & \texttt{brane representation: R\_\allowbreak{}A R\_\allowbreak{}B + R\_\allowbreak{}B R\_\allowbreak{}A == 2 eta17\_\allowbreak{}AB ID16 for R = (gamma\_\allowbreak{}mu (x) 1\_\allowbreak{}4, gamma5 (x) rho\_\allowbreak{}k)} & True & 18. \\
13 & \texttt{intertwiner: Gamma17[A].U == U.R\_\allowbreak{}A for A = 0..7, and U is unitary} & True & 18. \\
14 & \texttt{chi17 == U.(gamma5 (x) rho4567).U\textasciicircum{}dag; rho4567 Hermitian with eigenvalues \{-1, -1, 1, 1\}} & True & 18. \\
15 & \texttt{bilinears at the brane: psibar psi' == Sum\_\allowbreak{}a psibar\_\allowbreak{}a psi'\_\allowbreak{}a and psibar Gamma\textasciicircum{}mu psi' == Sum\_\allowbreak{}a psibar\_\allowbreak{}a gamma\textasciicircum{}mu psi'\_\allowbreak{}a (mu = 0..3)} & True & 18. \\
16 & \texttt{KK operator: X.X == -|q|\textasciicircum{}2/\allowbreak{}b0\textasciicircum{}2 ID16 and X anticommutes with gamma\_\allowbreak{}mu (x) 1\_\allowbreak{}4} & True & 18.1 \\
17 & \texttt{KK masses: Det[gamma.p (x) 1 + X - m] vanishes on p0\textasciicircum{}2 - |p|\textasciicircum{}2 == m\textasciicircum{}2 + |q|\textasciicircum{}2/\allowbreak{}b0\textasciicircum{}2 (40-digit arithmetic) and not off shell} & True & 18.1 \\
18 & \texttt{portal mixing: 1 - N.N\textasciicircum{}dag == mu.mu\textasciicircum{}dag/\allowbreak{}m\textasciicircum{}2 to second order (non-unitary light mixing)} & True & 18.2 \\
19 & \texttt{vacuum: P(a->b) doubly stochastic; equals delta - 4 Sum Re X sin\textasciicircum{}2 Delta + 2 Sum Im X sin 2Delta; CPT P(a->b) == P(bbar->abar)} & True & 19. \\
20 & \texttt{V\_\allowbreak{}CC(rho Y\_\allowbreak{}e = 1 g/\allowbreak{}cm\textasciicircum{}3) == 7.63 x 10\textasciicircum{}-14 eV; matter evolution doubly stochastic; a flavour-universal potential has no effect} & True & 19. \\
21 & \texttt{two-flavour MSW: sin\textasciicircum{}2 theta\_\allowbreak{}m from diagonalization == formula, and sin\textasciicircum{}2 2theta\_\allowbreak{}m == 1 at 2 E V == dm\textasciicircum{}2 cos 2theta} & True & 19. \\
22 & \texttt{Sun: numerical P\_\allowbreak{}ee == adiabatic MSW formula to 2 x 10\textasciicircum{}-3 at 0.3, 1, 5 and 10 MeV; P\_\allowbreak{}ee(0.3 MeV) within 0.03 of the vacuum average; P\_\allowbreak{}ee falls with energy} & True & 19. \\
23 & \texttt{cosmic neutrino background: T\_\allowbreak{}nu == 1.945 K and n == 112 cm\textasciicircum{}-3 per flavour} & True & 19. \\
24 & \texttt{dirac16complex portal: non-unitary oscillations, Sum\_\allowbreak{}b P(a->b) < 1 (zero-distance effect); decoupled limit == three-flavour vacuum result} & True & 19. \\*
\end{longtable}}
\textbf{The independent Python program} tools/ref\_sm.py — its boolean checks:

{\footnotesize
\begin{longtable}{>{\raggedright\arraybackslash}p{\dimexpr 0.0550\linewidth-2\tabcolsep\relax}>{\raggedright\arraybackslash}p{\dimexpr 0.8250\linewidth-2\tabcolsep\relax}>{\raggedright\arraybackslash}p{\dimexpr 0.1200\linewidth-2\tabcolsep\relax}}
\toprule
\textbf{\#} & \textbf{check} & \textbf{result} \\
\midrule\endfirsthead
\toprule
\textbf{\#} & \textbf{check} & \textbf{result} \\
\midrule\endhead
\endfoot
\bottomrule\endlastfoot
1 & \texttt{anomaly cancellation in each generation (with nu\_\allowbreak{}R): SU(3)\textasciicircum{}2U(1), SU(2)\textasciicircum{}2U(1), U(1)\textasciicircum{}3, grav\textasciicircum{}2U(1), SU(3)\textasciicircum{}3 all 0} & True \\
2 & \texttt{Witten SU(2) anomaly: number of doublets per generation (3 colours of Q + 1 L) = 4, even} & True \\
3 & \texttt{electric charges Q = T3 + Y: Q -> (2/\allowbreak{}3, -1/\allowbreak{}3), u\_\allowbreak{}R 2/\allowbreak{}3, d\_\allowbreak{}R -1/\allowbreak{}3, L -> (0, -1), e\_\allowbreak{}R -1, nu\_\allowbreak{}R 0} & True \\
4 & \texttt{effective number of relativistic degrees of freedom g* = 28 + (7/\allowbreak{}8) 90 = 106.75 (106.75 + 5.25 if nu\_\allowbreak{}R were thermal)} & True \\
5 & \texttt{gauge-boson mass matrix from |D<Phi>|\textasciicircum{}2: eigenvalues g\textasciicircum{}2v\textasciicircum{}2/\allowbreak{}4 (twice, W+-), (g\textasciicircum{}2+g'\textasciicircum{}2)v\textasciicircum{}2/\allowbreak{}4 (Z), 0 (photon)} & True \\
6 & \texttt{tree-level relations: m\_\allowbreak{}W = g v/\allowbreak{}2, m\_\allowbreak{}Z = sqrt(g\textasciicircum{}2+g'\textasciicircum{}2) v/\allowbreak{}2, rho = m\_\allowbreak{}W\textasciicircum{}2/\allowbreak{}(m\_\allowbreak{}Z\textasciicircum{}2 cos\textasciicircum{}2) = 1, m\_\allowbreak{}H\textasciicircum{}2 = 2 lambda v\textasciicircum{}2} & True \\
7 & \texttt{CKM unitarity V V\textasciicircum{}dag = 1} & True \\
8 & \texttt{J = c12 c23 c13\textasciicircum{}2 s12 s23 s13 sin(delta) = 3.12e-5 (PDG 2024: 3.12 +0.13 -0.12 e-5)} & True \\
9 & \texttt{Wolfenstein parameters agree with the PDG fit (A 0.826, rhobar 0.1591, etabar 0.3523) within their quoted uncertainties} & True \\
10 & \texttt{Yukawa matrices Y\_\allowbreak{}u = sqrt2 diag(m\_\allowbreak{}u,m\_\allowbreak{}c,m\_\allowbreak{}t)/\allowbreak{}v, Y\_\allowbreak{}d = sqrt2 V diag(m\_\allowbreak{}d,m\_\allowbreak{}s,m\_\allowbreak{}b)/\allowbreak{}v: singular values x v/\allowbreak{}sqrt2 = masses, |U\_\allowbreak{}uL\textasciicircum{}dag U\_\allowbreak{}dL| = |V|} & True \\
11 & \texttt{PMNS unitarity} & True \\
12 & \texttt{neutrino Yukawa matrix Y\_\allowbreak{}nu = sqrt2 U diag(m1,m2,m3)/\allowbreak{}v: masses and |U| recovered (singular value decomposition)} & True \\
13 & \texttt{vacuum: probabilities P(a -> b) form a doubly stochastic matrix (rows and columns sum to 1)} & True \\
14 & \texttt{vacuum: amplitude result = delta\_\allowbreak{}ab - 4 sum\_\allowbreak{}\{i>j\} Re(X\_\allowbreak{}ij) sin\textasciicircum{}2 Delta\_\allowbreak{}ij + 2 sum\_\allowbreak{}\{i>j\} Im(X\_\allowbreak{}ij) sin 2Delta\_\allowbreak{}ij, X\_\allowbreak{}ij = U\_\allowbreak{}ai* U\_\allowbreak{}bi U\_\allowbreak{}aj U\_\allowbreak{}bj*} & True \\
15 & \texttt{CPT: P(a->b) = P(bbar->abar); CP asymmetry P(mu->e) - P(mubar->ebar) proportional to J (vanishes for delta = 0, pi)} & True \\
16 & \texttt{V\_\allowbreak{}CC(rho Y\_\allowbreak{}e = 1 g/\allowbreak{}cm\textasciicircum{}3) = 7.63e-14 eV} & True \\
17 & \texttt{matter (constant density): unitary evolution, probabilities doubly stochastic} & True \\
18 & \texttt{matter: a flavour-universal potential has no effect (only differences of potentials matter)} & True \\
19 & \texttt{two-flavour MSW: sin\textasciicircum{}2 theta\_\allowbreak{}m from diagonalization = formula; maximal mixing (sin\textasciicircum{}2 2theta\_\allowbreak{}m = 1) at 2 E V = dm\textasciicircum{}2 cos 2theta} & True \\
20 & \texttt{Sun: numerical P\_\allowbreak{}ee = adiabatic MSW formula to 2e-3 at E = 0.3, 1, 5, 10 MeV} & True \\
21 & \texttt{Sun: P\_\allowbreak{}ee falls from near the vacuum average c13\textasciicircum{}4 (1 - sin\textasciicircum{}2(2 theta12)/\allowbreak{}2) + s13\textasciicircum{}4 (within 0.03 at 0.3 MeV) to about 0.31 at 10 MeV} & True \\
22 & \texttt{CNB: T\_\allowbreak{}nu = 1.945 K and n = 112 cm\textasciicircum{}-3 per flavour} & True \\
23 & \texttt{brane frame metric diag(+1,-1,-1,-1) = eta17 on x0..x3 = squares of the quaternion basis 1, e\_\allowbreak{}1, e\_\allowbreak{}2, e\_\allowbreak{}3 (Re(e\_\allowbreak{}a e\_\allowbreak{}b)); H closed} & True \\
24 & \texttt{brane representation: gamma\_\allowbreak{}mu (x) 1\_\allowbreak{}4 (mu = 0..3) and gamma5 (x) rho\_\allowbreak{}k (k = 4..7) satisfy \{R\_\allowbreak{}A, R\_\allowbreak{}B\} = 2 eta\_\allowbreak{}AB} & True \\
25 & \texttt{intertwiner U (Schur average over the 256 Clifford products): Gamma17[A] U = U R\_\allowbreak{}A for A = 0..7, U unitary} & True \\
26 & \texttt{8D chirality at the brane: chi17 = U (gamma5 (x) rho4567) U\textasciicircum{}dag, rho4567 Hermitian with eigenvalues +1, +1, -1, -1} & True \\
27 & \texttt{bilinears at the brane: Psibar Psi' = sum\_\allowbreak{}a psibar\_\allowbreak{}a psi'\_\allowbreak{}a and Psibar Gamma\textasciicircum{}mu Psi' = sum\_\allowbreak{}a psibar\_\allowbreak{}a gamma\textasciicircum{}mu psi'\_\allowbreak{}a (mu = 0..3)} & True \\
28 & \texttt{Kaluza-Klein: the 16x16 operator gamma.p (x) 1 + gamma5 (x) rho.q/\allowbreak{}b0 - m is singular exactly on p0\textasciicircum{}2 - |p|\textasciicircum{}2 = m\textasciicircum{}2 + |q|\textasciicircum{}2/\allowbreak{}b0\textasciicircum{}2} & True \\
29 & \texttt{KK operator: X\textasciicircum{}2 = -|q|\textasciicircum{}2/\allowbreak{}b0\textasciicircum{}2 and X anticommutes with gamma\textasciicircum{}mu (x) 1} & True \\
30 & \texttt{non-unitary oscillations: sum\_\allowbreak{}b P(a->b) = (N N\textasciicircum{}dag)\_\allowbreak{}aa\textasciicircum{}2-type deficit < 1 (zero-distance effect), complemented by the heavy states} & True \\
31 & \texttt{portal decoupling: for mu -> 0 the light block becomes the PMNS matrix and the oscillation probabilities the three-flavour vacuum ones (section 19)} & True \\
32 & \texttt{age 13.79 Gyr (Planck 2018: 13.787 +- 0.020)} & True \\
33 & \texttt{ADD relation: 3 Ha\textasciicircum{}2 = kappa8 rho4/\allowbreak{}V\_\allowbreak{}perp = 8 pi G4 rho4 with G4 = G8/\allowbreak{}V\_\allowbreak{}perp} & True \\*
\end{longtable}}
\textbf{Numbers printed by both programs} (relative differences; the Python program prints $|V|$ with 6 and $|U|$ with 5 decimals):

{\small
\begin{longtable}{>{\raggedright\arraybackslash}p{\dimexpr 0.6000\linewidth-2\tabcolsep\relax}>{\raggedright\arraybackslash}p{\dimexpr 0.1500\linewidth-2\tabcolsep\relax}>{\raggedright\arraybackslash}p{\dimexpr 0.2500\linewidth-2\tabcolsep\relax}}
\toprule
\textbf{quantity} & \textbf{values} & \textbf{max. relative difference} \\
\midrule\endfirsthead
\toprule
\textbf{quantity} & \textbf{values} & \textbf{max. relative difference} \\
\midrule\endhead
\endfoot
\bottomrule\endlastfoot
$v,g,g',\sin^2\theta_W,e,1/\alpha,\lambda$ & 7 & $2.1\times10^{-16}$ \\
$|V_{\rm CKM}|$ & 9 & $2.3\times10^{-5}$ \\
$J_{\rm CKM}$ & 1 & $2.2\times10^{-16}$ \\
$\bar\rho,\bar\eta$ & 2 & $3.5\times10^{-16}$ \\
$|U_{\rm PMNS}|$ & 9 & $3.2\times10^{-5}$ \\
$J_{\rm PMNS}$ & 1 & $1.2\times10^{-16}$ \\
$m_1,m_2,m_3$ (eV) & 3 & 0 \\
$y_t$ & 1 & $1.1\times10^{-16}$ \\
age (Gyr) & 1 & $1.0\times10^{-15}$ \\
$V_\perp$, $V_\perp^{1/4}$, lowest KK mass & 3 & $2.2\times10^{-16}$ \\
light (eV) and heavy (GeV) masses with the portal & 7 & $9.9\times10^{-7}$ \\
$(1-NN^\dagger)_{ee}$ & 1 & $3.3\times10^{-7}$ \\
DUNE vacuum $P(\mu\to e)$, $P(\bar\mu\to\bar e)$, $P(\mu\to\mu)$ & 3 & $4.3\times10^{-16}$ \\
DUNE matter $P(\mu\to e)$, $P(\bar\mu\to\bar e)$ & 2 & $1.2\times10^{-15}$ \\
vector portal: $P(\mu\to e)$ with, without & 2 & $9.0\times10^{-16}$ \\
Sun $P_{ee}$ ($E$, numerical, adiabatic) & 12 & $3.8\times10^{-6}$ \\
supernova resonance densities H, L & 2 & $1.9\times10^{-16}$ \\
$T_\nu$, $n_\nu$ & 2 & $1.3\times10^{-16}$ \\*
\end{longtable}}
\textbf{SHA-256 of the files} (at the time this document was generated):

\begin{Verbatim}[fontsize=\small,frame=single,framesep=2mm]
7+1.nb                             41d252f9852f210c228f45e1ac956d2d7ef26faeb0062a5880280e1aae0ef
    ...fc7
tools/make_7p1_nb.py               996d26c9552cd44902560af16a1683fa71457261f18fefaeddae87bd99a2b
    ...144
tools/nb_part_gravity.py           c4caf2ad5c932e8f603341edade5f4b98fa784feb05906c051eab7c884cd6
    ...8f8
tools/nb_part_ks_gravity.py        a656c0ad040b4576b75b17e029545109e531c8e656ad99bfdb5cba946e967
    ...772
tools/nb_part_sm.py                43e92beed57453a4b57931de6d603ac6c218efba6e4005133b69b8054a975
    ...c5b
tools/nb_part_coupled.py           e630521d64132e71afc861114f699eb06454557cf14a57ad8ec1eef96247a
    ...ae0
tools/nb_part_dark.py              10ad8fb8ae25fa07bbee758c653995acd4e2c80733059a8b3648bd1593e46
    ...db0
tools/nb_part_torus_brane.py       fb904ff8d85b904c55b3b2a8f60fb0cb047f1ccf1b98db3364092f2bae0b2
    ...162
tools/nb_part_radion.py            53c5604707a3e22228c840ee0a2d2125616f9410f75a5535041dfaf18f497
    ...c00
tools/lovelock71.py                d93ecbd0f82adb972218d59dfe32999a4af944945ca9758607f3a633e1af9
    ...5a1
tools/ref_gravity.py               f07cf8adaff585af89c87d19532f1aaa7db606ad88ae90483b2d1f41341ed
    ...f36
tools/ref_ks_gravity.py            b6488ef24b1c16aacd8ff6cf6a5a1ba3f316e082449ad97480022d1a26383
    ...27f
tools/ref_sm.py                    f3f80806893f42bbf166e1f1de896ffd611dbfe847685da8e3ee367e9cde6
    ...aba
tools/ref_coupled.py               695dbd9375753653472d9c37502f0067c4a64fe01a64ba8c05e8b2c5da85a
    ...a58
tools/ref_dark.py                  1600e6b57ab0383b4f2b664e2796d3864c9afa77a9ed59f007f66e1139453
    ...ff5
tools/ref_torus_brane.py           1f623656ab0f5c8ee43cb310f7bfd580cbba5a7385da2963adaae4d2f8bda
    ...af7
tools/ref_radion.py                e9d47176b9f2d2232a84b3b82188bd3b0a19cde3dbe1ba013c1fe5abff634
    ...1ce
tools/ref71.py                     1a7a06eb0add3f459611776dc9c474a5bb959d947a9db949db4612b1ab943
    ...4e0
tools/ref44.py                     ab100c3ebcc055fbe59c132a02f80ecb17630e934e99d59f88e919ec8064f
    ...af5
\end{Verbatim}
\textbf{Output of tools/ref\_sm.py (independent checks of sections 17-19; exit code 0)} (complete):

\begin{Verbatim}[fontsize=\small,frame=single,framesep=2mm]
anomaly cancellation in each generation (with nu_R): SU(3)^2U(1), SU(2)^2U(1), U(1)^3,
    ...grav^2U(1), SU(3)^3 all 0 : True
Witten SU(2) anomaly: number of doublets per generation (3 colours of Q + 1 L) = 4, even : True
electric charges Q = T3 + Y: Q -> (2/3, -1/3), u_R 2/3, d_R -1/3, L -> (0, -1), e_R -1, nu_R 0
    ...: True
effective number of relativistic degrees of freedom g* = 28 + (7/8) 90 = 106.75 (106.75 + 5.25
    ...if nu_R were thermal) : True
gauge-boson mass matrix from |D<Phi>|^2: eigenvalues g^2v^2/4 (twice, W+-), (g^2+g'^2)v^2/4
    ...(Z), 0 (photon) : True
NUM v = (sqrt2 G_F)^(-1/2) [GeV], g, g', sin^2 theta_W (on shell), e, 1/alpha(tree), lambda =
    ...m_H^2/(2v^2) : (246.21964023926205, 0.6528252573344827, 0.34994623401223707,
    ...0.22320949284344327, 0.3084276316025824, 132.10014858375337, 0.1292805756612353)
tree-level relations: m_W = g v/2, m_Z = sqrt(g^2+g'^2) v/2, rho = m_W^2/(m_Z^2 cos^2) = 1,
    ...m_H^2 = 2 lambda v^2 : True
CKM unitarity V V^dag = 1 : True
NUM CKM |V| (rows ud us ub / cd cs cb / td ts tb) : [[0.97435, 0.225008, 0.003732], [0.224876,
    ...0.973489, 0.04183], [0.008583, 0.041109, 0.999118]]
NUM Jarlskog J = Im(V_us V_cb V_ub* V_cs*) : 3.1169935287554706e-05
J = c12 c23 c13^2 s12 s23 s13 sin(delta) = 3.12e-5 (PDG 2024: 3.12 +0.13 -0.12 e-5) : True
NUM Wolfenstein lambda = s12, A = s23/s12^2, rhobar + i etabar = -V_ud V_ub*/(V_cd V_cb*) :
    ...(0.22501, 0.8261981634695225, np.float64(0.15918443594247167),
    ...np.float64(0.35227461890586964))
Wolfenstein parameters agree with the PDG fit (A 0.826, rhobar 0.1591, etabar 0.3523) within
    ...their quoted uncertainties : True
Yukawa matrices Y_u = sqrt2 diag(m_u,m_c,m_t)/v, Y_d = sqrt2 V diag(m_d,m_s,m_b)/v: singular
    ...values x v/sqrt2 = masses, |U_uL^dag U_dL| = |V| : True
NUM top Yukawa y_t = sqrt2 m_t/v : 0.9911341438317451
PMNS unitarity : True
NUM PMNS |U| and J_PMNS : ([[0.82328, 0.54796, 0.14816], [0.27491, 0.61298, 0.74073], [0.49663,
    ...0.56921, 0.65526]], np.float64(0.0017359130944756358))
NUM Dirac neutrino masses (lightest m1 = 0, eV) and their sum : (0.0, 0.008654478609367522,
    ...0.05033885179461287, 0.058993330403980394)
neutrino Yukawa matrix Y_nu = sqrt2 U diag(m1,m2,m3)/v: masses and |U| recovered (singular
    ...value decomposition) : True
NUM largest neutrino Yukawa sqrt2 m3/v : 2.891316340688847e-13
vacuum: probabilities P(a -> b) form a doubly stochastic matrix (rows and columns sum to 1) :
    ...True
vacuum: amplitude result = delta_ab - 4 sum_{i>j} Re(X_ij) sin^2 Delta_ij + 2 sum_{i>j}
    ...Im(X_ij) sin 2Delta_ij, X_ij = U_ai* U_bi U_aj U_bj* : True
NUM DUNE vacuum (L = 1300 km, E = 2.5 GeV): P(mu->e), P(mubar->ebar), P(mu->mu) :
    ...(np.float64(0.04883685296656589), np.float64(0.05019851034574761),
    ...np.float64(0.012813139189347619))
CPT: P(a->b) = P(bbar->abar); CP asymmetry P(mu->e) - P(mubar->ebar) proportional to J
    ...(vanishes for delta = 0, pi) : True
NUM V_CC = sqrt2 G_F n_e for rho Y_e = 1 g/cm^3 [eV] (textbook 7.63e-14 eV) :
    ...7.632466872241316e-14
V_CC(rho Y_e = 1 g/cm^3) = 7.63e-14 eV : True
matter (constant density): unitary evolution, probabilities doubly stochastic : True
NUM DUNE in matter (rho = 2.848 g/cm^3, Y_e = 0.5): P(mu->e), P(mubar->ebar) :
    ...(np.float64(0.07232693782692827), np.float64(0.029191799469459462))
matter: a flavour-universal potential has no effect (only differences of potentials matter) :
    ...True
two-flavour MSW: sin^2 theta_m from diagonalization = formula; maximal mixing (sin^2 2theta_m =
    ...1) at 2 E V = dm^2 cos 2theta : True
NUM Sun: P_ee (numerical through the solar density profile, adiabatic formula) for E = 0.3, 1,
    ...5, 10 MeV : {0.3: (np.float64(0.5251624291065616), 0.5251645055880894), 1.0:
    ...(np.float64(0.4584114174121251), 0.4584122490661078), 5.0:
    ...(np.float64(0.3103377102301538), 0.3103335111666083), 10.0:
    ...(np.float64(0.2979217906988696), 0.29791748878645763)}
Sun: numerical P_ee = adiabatic MSW formula to 2e-3 at E = 0.3, 1, 5, 10 MeV : True
Sun: P_ee falls from near the vacuum average c13^4 (1 - sin^2(2 theta12)/2) + s13^4 (within
    ...0.03 at 0.3 MeV) to about 0.31 at 10 MeV : True
NUM supernova MSW resonance densities for E = 20 MeV, Y_e = 0.5: H (dm31, theta13), L (dm21,
    ...theta12) [g/cm^3] : (1587.1391520947036, 18.939748107619376)
NUM cosmic neutrino background: T_nu = (4/11)^(1/3) T_gamma [K], number density per flavour
    ...[cm^-3] : (1.9453688391750839, 112.01641316574336)
CNB: T_nu = 1.945 K and n = 112 cm^-3 per flavour : True
brane frame metric diag(+1,-1,-1,-1) = eta17 on x0..x3 = squares of the quaternion basis 1,
    ...e_1, e_2, e_3 (Re(e_a e_b)); H closed : True
brane representation: gamma_mu (x) 1_4 (mu = 0..3) and gamma5 (x) rho_k (k = 4..7) satisfy
    ...{R_A, R_B} = 2 eta_AB : True
intertwiner U (Schur average over the 256 Clifford products): Gamma17[A] U = U R_A for A =
    ...0..7, U unitary : True
8D chirality at the brane: chi17 = U (gamma5 (x) rho4567) U^dag, rho4567 Hermitian with
    ...eigenvalues +1, +1, -1, -1 : True
bilinears at the brane: Psibar Psi' = sum_a psibar_a psi'_a and Psibar Gamma^mu Psi' = sum_a
    ...psibar_a gamma^mu psi'_a (mu = 0..3) : True
Kaluza-Klein: the 16x16 operator gamma.p (x) 1 + gamma5 (x) rho.q/b0 - m is singular exactly on
    ...p0^2 - |p|^2 = m^2 + |q|^2/b0^2 : True
KK operator: X^2 = -|q|^2/b0^2 and X anticommutes with gamma^mu (x) 1 : True
portal mixing: light masses (eV) and heavy masses (GeV) of the 7x7 Dirac mass matrix : ([0.0,
    ...0.008654465673872413, 0.05033886016084743], [1.0, 1.0, 1.00000000025,
    ...1.0000000007499998])
portal mixing: N N^dag = 1 - mu mu^dag/m^2 to second order (non-unitary light mixing), (1 - N
    ...N^dag)_ee : (True, np.float64(9.99999860695766e-10))
non-unitary oscillations: sum_b P(a->b) = (N N^dag)_aa^2-type deficit < 1 (zero-distance
    ...effect), complemented by the heavy states : True
portal decoupling: for mu -> 0 the light block becomes the PMNS matrix and the oscillation
    ...probabilities the three-flavour vacuum ones (section 19) : True
NUM vector portal (V_P,e = 0.2 V_CC): P(mu->e) at DUNE with and without the dirac16complex
    ...potential : (np.float64(0.07704998293140915), np.float64(0.07232693782692827))
NUM age of the universe from the brane Friedmann equation (Planck 2018: H0 = 67.4, Omega_m =
    ...0.315) [Gyr] : 13.790728126409226
age 13.79 Gyr (Planck 2018: 13.787 +- 0.020) : True
ADD relation: 3 Ha^2 = kappa8 rho4/V_perp = 8 pi G4 rho4 with G4 = G8/V_perp : True
NUM ADD example (reduced 8D Planck mass 10 TeV): V_perp [GeV^-4], V_perp^(1/4) [m], lowest KK
    ...mass 2 pi/V_perp^(1/4) [MeV] : (5929225000000.0, 3.079186208511254e-13,
    ...4.026524867292262)
boolean checks: 33   failed: 0
VERDICT: ALL CHECKS PASSED
\end{Verbatim}
\end{document}
